\chapter{Early User Personas}
\label{app:personas}

The group created two personas to explore different needs and expectations for Omnisound. One represents someone who regularly manages different types of audio, while the other represents someone who mainly wants quick access to music. The personas are design profiles rather than participants in the later usability tests.

\section{Persona A: George Flambae'}

George is a 29-year-old university student studying music theory. He is interested in audio and is comfortable using phone applications and computers.

His main goals are to find specific content, organise saved items and continue listening from where he previously stopped. His frustrations include having to search through separate libraries, forgetting where content was saved and losing track of unfinished listening.

For Omnisound, this persona highlights the importance of a shared Library, Search with clear content types, visible listening progress and an easy way to resume unfinished audio.

\section{Persona B: Hyleen Trisha}

Hyleen is a 24-year-old university student studying for a master's degree in computer engineering. She is very comfortable using computers and prefers finding the music she wants without unnecessary steps.

Her main goals are to reach music quickly and find familiar playlists. A crowded Home screen and unrelated podcasts or audiobooks competing for attention are her main frustrations.

For Omnisound, this persona highlights the importance of clear category navigation, recognisable music content and a Home screen that does not overwhelm users with unrelated information.

\section{Design Considerations}

The two personas illustrate why Omnisound needs to support different listening habits. A shared Library and Continue Listening help users manage several types of audio, while separate categories make it possible to focus on music, podcasts or audiobooks.

The prototype addresses these needs through Search, Library, clear categories and listening progress. Customising or hiding individual Home sections was not implemented in the tested versions.
